% ---
% id: reductions/domination/rem:rotated-projective
% kind: rem
% title: "Rotated projective elements and the noise floor"
% summary: >-
%   A projective element on a direction with no Gaussian-rational vector
%   dominates no non-zero lattice element, and completeness then forces a
%   weight of one, so no domination link runs from that block. A noise floor
%   in the model can remove it; after another channel it applies anew.
% prominence: standard
% status: accepted
% uses: [def:domination-link, foundations/def:domination, foundations/lem:domination, foundations/lem:rational-obstruction, foundations/lem:absorbing-domination]
% mentions: [foundations/rem:noise-floor, receiver/detector-imperfections/prop:noise-floor]
% assumes: []
% refs: []
% provenance: newly-derived
% lean:
% notes: "The remark follows foundations/lem:rational-obstruction's lead-in (after another channel the obstruction applies to the image of the reference in the new basis), foundations/rem:noise-floor (the factor of outcomes and dimension in the order of the weight) and receiver/detector-imperfections/prop:noise-floor (the conditions under which the noise-floor remedy works; loss alone or dark counts alone need not suffice). No requirement uses the remark: no run relates a snap to its model by domination."
% ---

The remark is about domination only, the relation the link of
Definition~\ref{def:domination-link} carries, whose channel acts on each
sector as the identity followed by a branching. A reduction that first
applies another channel to the reference changes what is compared: the
obstruction then applies to the image of the reference, compared with a
lattice family written in the new basis, where the supports may have
rational directions that they lacked. Such reductions are planned beyond
these contracts (\repodoc{spec/future-scope.md}).

\begin{remark}\label{rem:rotated-projective}
Let a reference family $G$, complete at $1$ on a block, have an element
$G_j=\dyad{\psi}$ that is a rank-one projector on a direction containing no
non-zero vector of $\mathbb{Q}(i)^{d}$, as happens for an ideal analyser
whose analysed polarisation $\cos\alpha\ket{H}+\sin\alpha\ket{V}$ has an
irrational $\tan\alpha$ on its one-photon block, or for the
two-photon block of a diagonal-basis analyser written in the orthonormal
occupation basis. Any family $F$ with Gaussian-rational entries that $G$
dominates at a weight below one has
$\mathrm{supp}\,F_j\subseteq\mathrm{supp}\,G_j$ by
\xref{foundations/lem:domination}\,(ii), so $F_j=\lambda\dyad{\psi}$ with
$\lambda\ge0$, and a non-zero such $F_j$ would have a column that is a
non-zero Gaussian-rational vector on the direction $\psi$; hence $F_j=0$:
$G_j$ dominates no non-zero lattice element. Completeness then forces the
weight up to one: \xref{foundations/lem:rational-obstruction}\,(v) gives
$\vartheta\ge\lmax(G_j)/1=1$, so no weight below one exists, at any lattice
exponent. No domination link runs from the reference on a block that holds
such an element, and no lattice family is related to it by domination there;
this is the obstruction \xref{foundations/rem:noise-floor} explains. The
remedy within domination is in the model and not in the lattice. With the
device's loss and dark counts composed into the reference, every element is
positive definite on every block under the conditions that the receiver
module states (\xref{receiver/detector-imperfections/prop:noise-floor}:
loss in the receiver and a dark-count probability strictly between zero and
one in every port); loss alone or dark counts alone need not suffice. Where
those conditions hold, the obstruction disappears: a lattice family whose
elements but one lie between zero and their reference elements, the last
taking what completeness leaves, is dominated at a weight of at most
$\delta/(\lambda+\delta)$, with $\delta$ the norm of the residue that last
element absorbs and $\lambda$ the least eigenvalue of the reference's element
at the same outcome, and not of the lattice element, which
the residue raises above it
(\xref{foundations/lem:absorbing-domination}\,(iv);
\xref{foundations/rem:noise-floor}).
After another channel the obstruction is to be read in the new basis, where
it may not arise: an ideal analyser at an angle with irrational
$\tan\alpha$ is the
image, under one unitary applied before every outcome, of the same analyser
at an angle where its elements are lattice values, and the isometry from the
two-photon block into $\mathbb{C}^{2}\otimes\mathbb{C}^{2}$ carries the
diagonal-basis element above to a projector on a rational vector; neither
route is a claim of this contract. The \tref{emission normalizer} does not meet the
obstruction in the form above: the datum it snaps has Gaussian-rational
entries, so every support it meets is spanned by Gaussian-rational vectors
(\xref{foundations/lem:rational-obstruction}\,(ii)).
\end{remark}
